One mitigation is carried through the design and pre-registered as $M1$: increase patch overlap $O$ by reducing the stride $S$ at fixed patch size $P$. Its geometric basis lies in the stride branch, whose candidates are in $\mathcal{F}_{S}$. $M1$ asks whether this reduction also decreases the observed recovery deficit. Model~A$'$ of \autoref{sec:three} tests the claim through the overlap slope $\delta_O$, read on the contrast $d$ that Model~A is fitted to, where a mitigation is a positive slope. The pre-registered rule has two legs: $\Pr(\delta_O>0\mid D)\geq0.95$, and the configuration level carrying the overlap term must win the leave-one-out comparison against the same level without it. Both must hold.

Part of the answer does not wait on the fit. In the analysed band the stride candidates satisfy $f=cf_s/S$ with $c<S/2$, so a geometry has $\lceil S/2\rceil-1$ of them; a smaller $S$ raises the fundamental $f_s/S$, widens the spacing and lowers that count, and halving $S$ retains only the even-indexed sites. The patch fundamental $f_s/P$ is untouched by any of this, so $\mathcal{F}_{P}$ is invariant under the mitigation by construction. Overlap is therefore a partial geometric adjustment at best: it displaces stride-derived candidates and cannot remove patch-derived ones.
